\documentclass{article}

\usepackage{fancyhdr}
\usepackage{amsmath}
\usepackage{amsthm}
\usepackage{amsfonts}
\usepackage{tikz}
\usepackage[plain]{algorithm}
\usepackage{algpseudocode}

\usetikzlibrary{automata,positioning}

%
% Basic Document Settings
%

\topmargin=-0.45in
\evensidemargin=0in
\oddsidemargin=0in
\textwidth=6.5in
\textheight=9.0in
\headsep=0.25in

\linespread{1.1}

\pagestyle{fancy}
\lhead{\hmwkAuthorName}
\chead{} % empty
\rhead{\hmwkHeaderTitle}
\cfoot{\thepage}

\renewcommand\headrulewidth{0.4pt}
\renewcommand\footrulewidth{0pt}

\setlength\parindent{0pt}

%
% Homework Details
%

\newcommand{\hmwkTitle}{Homework\ 4}
\newcommand{\hmwkDueDate}{Tuesday, February 24, 2026,\ 11{:}59 pm (Thai time)}
\newcommand{\hmwkClass}{Data Structure\ and Abstractions}
\newcommand{\hmwkClassTime}{Section 2}
\newcommand{\hmwkClassInstructor}{Kanat Tangwongsan}
\newcommand{\hmwkAuthorName}{\textbf{Jiraroj Wiruchpongsanon (6781617)}}

% What appears in the top-right header
\newcommand{\hmwkHeaderTitle}{Data Structure and Abstractions: Homework 4}

%
% Title Page
%

\title{
    \vspace{2in}
    \textmd{\textbf{\hmwkClass:\ \hmwkTitle}}\\
    \normalsize\vspace{0.1in}\small{Due\ on\ \hmwkDueDate}\\
    \vspace{0.1in}\large{\textit{\hmwkClassInstructor\ \hmwkClassTime}}
    \vspace{3in}
}

\author{\hmwkAuthorName}
\date{}

\renewcommand{\part}[1]{\textbf{\large Part \Alph{partCounter}}\stepcounter{partCounter}\\}

%
% Helper Commands
%

\newcounter{partCounter}
\newcommand{\solution}{\textbf{\large Solution}}

\newtheorem*{theorem}{Theorem}

\theoremstyle{remark}
\newtheorem*{remark}{Remark}

\begin{document}

\maketitle
\newpage

\section*{Task 1: Missing Tile (6 points)}

For this task, save your work in \texttt{smallman/MissingTile.java}. There are four ways an L-shaped triomino can be arranged, labeled by the missing corner: NE, SE, SW, or NW.

\begin{theorem}
Any $2^n \times 2^n$ grid with one painted cell can be tiled using L-shaped triominoes such that the entire grid is covered by triominoes but no triomino overlaps either another triomino or the painted cell.
\end{theorem}

\textbf{Subtask I:} Re-prove the theorem above, in your own words, by induction on $n$ so that you fully understand how it works.

\textbf{Subtask II:} Turn the proof into code. Inside \texttt{MissingTile.java}, implement
\begin{verbatim}
public static void tileGrid(Grid board)
\end{verbatim}
that tiles the provided board instance using the information obtained from the instance. The board instance implements the following \texttt{Grid} interface:
\begin{verbatim}
// The top-left corner is coordinate x=0 and y=0. The bottom-right corner is
// coordinate x=size()-1 and y=size()-1.
public interface Grid {
    // return the width/height of the square grid. coordinates in the grid are
    // numbers between 0 and size()-1, inclusive.
    int size();
    // return the x coordinate of the painted cell
    int getPaintedCellX();
    // return the y coordinate of the painted cell
    int getPaintedCellY();
    // install a triomino tile so that the missing corner is at coordinate (x, y)
    // and the orientation of the tile is given by the orientation parameter,
    // encoded as follows: missing NE = 0, missing SE = 1, missing SW = 2, and
    // missing NW = 3. The function returns true if the tiling is possible; it
    // returns false if the tiling failed (e.g., overlap with an existing tile or
    // touches the painted cell.
    boolean setTile(int x, int y, int orientation);

    // return a boolean value indicating whether the grid has been fully tiled
    boolean isFullyTiled();
}
\end{verbatim}

As an example, suppose you are given a board instance whose size is $4 \times 4$ and the painted cell is at coordinate $x = 3$ and $y = 1$. That is, \texttt{board.size()} returns 4, \texttt{board.getPaintedCellX()} returns 3, and \texttt{board.getPaintedCellY()} returns 1. The grid can be manually tiled by invoking
\begin{enumerate}
    \item \texttt{board.setTile(1, 1, 1)}
    \item \texttt{board.setTile(3, 1, 1)}
    \item \texttt{board.setTile(2, 1, 0)}
    \item \texttt{board.setTile(1, 2, 0)}
    \item \texttt{board.setTile(2, 2, 3)}
\end{enumerate}
Here, the argument \texttt{orientation} uses the encoding described in the interface specification above (e.g., \texttt{1} means the triomino is missing the SE corner).

\textbf{Performance Expectations:} Your code will be tested with grid sizes between $1 \times 1$ and $1024 \times 1024$. A call to \texttt{tileGrid} should return within 2 seconds.

\textbf{Test Driver:} The starter pack provides a Java program that implements the \texttt{Grid} interface. You should not need to modify this program. Keep in mind that the grader will use the original version of the driver (\texttt{TileDriver.java}).

\textbf{Getting Started Guide:} The inductive proof shows that to tile a $2^k \times 2^k$ grid, we recursively tile four $2^{k-1} \times 2^{k-1}$ subgrids of the big grid, each with a different painted cell. The key question is how to account for each subgrid when its painted cell potentially differs from the original grid. Consider the following design approaches (pick one):
\begin{itemize}
    \item \textbf{Idea \#1:} Write a helper method with parameters such as
\begin{verbatim}
static void tilingHelper(int size, int topX, int topY,
                         int paintedX, int paintedY, Grid bd)
\end{verbatim}
    that explicitly passes the coordinates describing the current subgrid and its painted cell.
    \item \textbf{Idea \#2:} Instead of passing many parameters, create and act on a view of the real object. Implement an inner class \texttt{GridView} that also implements \texttt{Grid}. Each time you subdivide the grid, return a \texttt{GridView} instance whose \texttt{size}, \texttt{getPaintedCellX}, and \texttt{getPaintedCellY} reflect the subgrid's information, and whose \texttt{setTile} method forwards the command to the real grid while compensating for offsets. You do not need \texttt{isFullyTiled} to work for the view. The class should at least support
\begin{verbatim}
public GridView(Grid parentGrid);
public GridView subgrid(int topX, int topY, int size,
                        int virtualPaintedX, int virtualPaintedY);
\end{verbatim}
    and any additional helpers you find useful.
\end{itemize}

\medskip
\solution

\section*{Task 4: Tail Sum of Squares (2 points)}

Consider the following snippet of Java code:
\begin{verbatim}
int sumHelper(int n , int a) {
    if (n==0) return a;
    else return sumHelper(n-1, a + n*n);
}
int sumSqr(int n) { return sumHelper(n, 0); }
\end{verbatim}

Your Task: Prove that for $n \ge 1$, \texttt{sumSqr(n)} returns $1^2 + 2^2 + 3^2 + \cdots + n^2$. To prove this, use induction to show that \texttt{sumHelper} computes the right thing. (Hint: Think about how we proved \texttt{fact\_helper} in class.)

\medskip
\solution

\section*{Task 5: Mysterious Function (2 points)}

Consider the following Python function \texttt{foo}, which takes as input an integer $n \ge 1$ and returns a tuple of length 2 of integers:
\begin{verbatim}
def foo(n):
    assert n>=1
    if n == 1:
        return (1, 2)
    else:
        p, q = foo(n-1)
        return (q + p*n*(n+1), q*n*(n+1))
\end{verbatim}

Prove that for $n \ge 1$, \texttt{foo(n)} returns a pair $(p, q)$ such that
\[
    \frac{p}{q} = 1 - \frac{1}{n + 1}.
\]
(Hint: Use induction on $n$.)

\medskip
\solution

\section*{Task 7: Midway Tower of Hanoi (4 points)}
For this task, save your work in \texttt{smallman/Midway.java}. 

Monks at a remote monastery are busy solving the Tower of Hanoi problem, a duty passed on for tens of generations. According to an old tale, when this group of monks finishes, $P$ will be shown to equal $NP$, and the world will come to an end.

In Tower of Hanoi, you are given $N$ disks labeled $0, 1, \ldots, N-1$ by their sizes. There are 3 pegs: Peg 0, Peg 1, and Peg 2. Initially, all disks are on Peg 0, neatly arranged from small (Disk 0) to large (Disk $N-1$), with the smallest on top. The goal is to transfer all disks to Peg 1. You can move exactly one disk at a time, and a larger disk can never be placed on top of a smaller disk.

Solving this task is straightforward with the following Python program, which prints the fewest number of moves---that is, the optimal solution:
\begin{verbatim}
def solve_hanoi(n, from_peg, to_peg, aux_peg):
    if n>0:
        solve_hanoi(n-1, from_peg, aux_peg, to_peg)
        print("Move disk", n-1, "from Peg", from_peg, "to Peg", to_peg)
        solve_hanoi(n-1, aux_peg, to_peg, from_peg)

solve_hanoi(n, 0, 1, 2)
\end{verbatim}
Following these steps uses $2^N - 1$ moves. Because this number is huge, the monks devised a puzzle: assume they strictly followed the program's instructions. Given an intermediate configuration, determine how many more steps they will need before completing the task. (The original handout includes a visual trace of all configurations generated by the program.)

\textbf{Subtask I:} Prove, using mathematical induction, that for any $n \ge 0$, \texttt{solve\_hanoi(n, ...)} generates exactly $2^n - 1$ lines of instructions.

\medskip
\solution

\begin{proof}

    we claim that

    \[
        P(n) \equiv \text{solve\_hanoi(n, ...) will print out for exactly } 2^n - 1 \text{ lines when } n \ge 0
    \]

    for the \textbf{base case} where $n_0 = 0$

    \[
        P(0) \equiv 2^0 - 1 = 0 \text{ as the function only works when } n > 0 \text{ so nothing got printed out}
    \]

    then, we assume P(k) for arbitrary $(k \ge n_0)$, we prove $P(k+1) = 2^{k+1} - 1$

    \vspace{1em}

    from the recurrence, if $n = 0$, the function does nothing, so nothing got printed out

    \[T(0) = 0\]

    and for $n > 0$, solve\_hanoi(n-1, ...) will print out $T(n-1)$ lines (moves the top n-1 disks to the left or right, so we can move disk n to the middle), then the mother function prints out (move disk n to the middle), and ended with another solve\_hanoi(n-1, ...) (with different parameter), this $T(n-1)$ lines of actually solving for $T(n-1)$ --- which the pile's position will iterate between left and right; thus

    \[T(n) = 2T(n-1) + 1)\]

    we use this recurrence relation to prove $P(k+1)$

    \begin{align*}
        P(k+1) &\equiv T(k+1) = 2T(k) + 1 && \quad \text{from the recurrence}\\
        &= 2(2^k-1) + 1 && \quad \text{induction hypothesis} \\
        &= 2^{k+1} - 2 + 1 = 2^{k+1} - 1
    \end{align*}

    since P(k+1) follows P(k), the \textbf{induction case} is satisfied, thus we can conclude the given predicate is true \end{proof}

\end{proof}
\end{document}
